% !TeX root = ../../tesis.tex

% =============================================================================
%  4.2.2 — ARQUITECTURA
% =============================================================================
\subsection{Arquitectura del Sistema}
\label{sec:vft-arquitectura}

\subsubsection{Tres Capas con Responsabilidades Delimitadas}
\label{subsec:vft-capas}


VFTModel está organizado bajo el patrón de \termino{arquitectura hexagonal} (
\termino{Ports and Adapters}). La idea central es que el núcleo matemático del
sistema —los algoritmos que calculan los indicadores topológicos— no debe
depender de cómo lleguen los datos ni de cómo se presenten los resultados. Esta
independencia se logra separando el código en tres capas concéntricas:

\begin{description}
    \item[\textbf{Capa de \api{}} (\texttt{src/api/})]
    Recibe las peticiones \textsc{http}, valida los datos de entrada mediante
    contratos de tipo estrictos y devuelve respuestas \textsc{json}. Es la única
    capa visible para quien consume el sistema desde el exterior.

    \item[\textbf{Capa de Dominio} (\texttt{src/core/})]
    Contiene los algoritmos topológicos, los modelos matemáticos y los servicios
    de construcción del grafo. Aquí reside toda la inteligencia analítica del
    modelo. No depende de cómo lleguen los datos ni de cómo se muestren.

    \item[\textbf{Capa de Infraestructura} (\texttt{src/infrastructure/})]
    Gestiona los clientes \textsc{http} para comunicarse con Apimetro, el acceso
    a archivos \geojson{} locales de respaldo y las colas de procesamiento
    asíncrono. Se encarga de obtener los datos, no de analizarlos. El cliente
    HTTP y la lógica de \termino{fallback} local se documentan en el
    Anexo~\ref{anx:vft-goclient}; la configuración de contenedores Docker y la
    automatización con \textit{Makefile} en el Anexo~\ref{anx:vft-infraestructura}.
\end{description}

La Figura~\ref{fig:vft-arquitectura} ilustra la jerarquía de las tres capas y el
sentido de las dependencias entre ellas.

\begin{figure}[H]
    \centering
    \resizebox{\linewidth}{!}{%
        % Diagrama: Arquitectura hexagonal de VFTModel (v4)
% Paleta: vftVerde (#059669, esmeralda) + vftGris (#6B7280, neutro)
% R=4.0 cm — zona segura (apotema=3.464) soporta text width=4.0 hasta y=−2.8

\definecolor{vftVerde}{HTML}{059669}
\definecolor{vftGris}{HTML}{6B7280}

\begin{tikzpicture}[
    adaptadorP/.style={
        rectangle, draw=vftVerde, fill=vftVerde!8,
        rounded corners=5pt, text width=3.6cm, align=left,
        minimum height=5.8cm, inner sep=11pt, font=\small
    },
    adaptador/.style={
        rectangle, draw=vftVerde!65, fill=vftVerde!4,
        rounded corners=4pt, text width=3.2cm, align=left,
        minimum height=2.8cm, inner sep=9pt, font=\small
    },
    adaptadorViz/.style={
        rectangle, draw=vftGris!55, fill=vftGris!5,
        rounded corners=5pt, text width=7.0cm, align=center,
        minimum height=2.6cm, inner sep=12pt, font=\small
    },
    actor/.style={
        rectangle, draw=vftGris!55, fill=vftGris!7,
        rounded corners=4pt, text width=2.1cm, align=center,
        minimum height=3.0cm, font=\small
    },
    flecha/.style={Stealth-Stealth, thick, color=vftVerde!65},
    flechaext/.style={Stealth-Stealth, thick, color=vftGris!50, dashed},
    flechaviz/.style={-Stealth, thick, color=vftGris!60, dashed}
]

%% ——— Hexágono central: Dominio (src/core/) ———
%% Flat-top, R = 4.0 cm
%% Vértices izq/der: (±4.0, 0)
%% Lados sup/inf: (±2.0, ±3.464)   [R·√3/2 = 4.0 · 0.866 = 3.464]
\draw[fill=vftVerde!10, draw=vftVerde, line width=1.5pt]
    ( 4.0,  0.000) --
    ( 2.0,  3.464) --
    (-2.0,  3.464) --
    (-4.0,  0.000) --
    (-2.0, -3.464) --
    ( 2.0, -3.464) -- cycle;

%% — Contenido del hexágono —
%% "Ports & Adapters" como tiny italic justo bajo src/core/, con holgura
\node[align=center, text width=4.0cm] at (0, 0) {
    \textbf{\large Dominio}\\[6pt]
    {\scriptsize\ttfamily src/core/}\\[4pt]
    {\tiny\itshape\textcolor{vftVerde!50!black}{Ports \& Adapters}}\\[10pt]
    \textcolor{vftVerde!45!black}{\rule{3.6cm}{0.3pt}}\\[8pt]
    {\small VFTGraphBuilder}\\[3pt]
    {\scriptsize \textit{STRICT} $\cdot$ \textit{REALISTIC}}\\[5pt]
    {\small Impedancia Vial}\\[8pt]
    \textcolor{vftVerde!45!black}{\rule{3.6cm}{0.3pt}}\\[8pt]
    {\scriptsize\textcolor{vftVerde!80!black}{Indicadores:}}\\[7pt]
    {\small $C$ — Cobertura Espacial}\\[6pt]
    {\small $k_{\text{in}}$ — Fuerza Capilar}
};

%% — Nodo invisible para posicionamiento y fit Docker —
\node[minimum width=8.0cm, minimum height=6.928cm,
      inner sep=0pt, draw=none, fill=none] (hexbox) at (0, 0) {};

%% ——— Adaptador primario (driving): Capa API ———
\node[adaptadorP, left=1.6cm of hexbox] (api) {
    \textbf{Capa \api{}}\\
    {\scriptsize\ttfamily src/api/}\\[7pt]
    FastAPI + Uvicorn\\
    Schemas Pydantic\\[7pt]
    \textcolor{vftVerde!70!black}{\scriptsize Endpoints \textsc{rest}:}\\[4pt]
    {\scriptsize\ttfamily /network/build-auto}\\[1pt]
    {\scriptsize\ttfamily /topological/*}\\[1pt]
    {\scriptsize\ttfamily /spatial-coverage}\\[1pt]
    {\scriptsize\ttfamily /geo-layers}\\[1pt]
    {\scriptsize (7~capas GeoJSON)}
};

%% ——— Adaptadores secundarios (driven): Infraestructura ———
%% yshift ±3.2 cm + height 2.8 cm → hueco 0.4 cm entre cajas
\node[adaptador, right=1.6cm of hexbox, yshift= 3.2cm] (goclient) {
    \textbf{Go Client}\\
    {\scriptsize\ttfamily go\_client/}\\[5pt]
    {\scriptsize Fetch estaciones,}\\
    {\scriptsize líneas, polígonos}\\[3pt]
    {\scriptsize $\to$ Apimetro \api{}}
};
\node[adaptador, right=1.6cm of hexbox, yshift= 0.0cm] (geojson) {
    \textbf{GeoJSON local}\\
    {\scriptsize\ttfamily geojson\_service/}\\[5pt]
    {\scriptsize Fallback: \texttt{map.geojson}}\\
    {\scriptsize sin conexión de red}
};
\node[adaptador, right=1.6cm of hexbox, yshift=-3.2cm] (cache) {
    \textbf{Persistencia}\\
    {\scriptsize\ttfamily persistence/}\\[5pt]
    {\scriptsize Redis (caché grafos)}\\
    {\scriptsize Celery (colas \textit{async})}
};

%% ——— Actores externos ———
\node[actor, left=1.3cm of api] (cliente) {
    Urbanista\\/ Analista
};
\node[actor, right=1.3cm of goclient] (apimetro) {
    \textbf{Apimetro}\\[3pt]
    {\scriptsize Go \textsc{rest}}
};

%% ——— Conector de Visualización (exterior al Docker, abajo) ———
\node[adaptadorViz, below=2.2cm of hexbox] (viz) {
    \textbf{Conectores de Visualización}\\[5pt]
    {\small Transport-GIS-MJG}\\[2pt]
    {\small Tableau $\cdot$ \textsc{qgis}}\\[5pt]
    {\scriptsize\textcolor{vftGris!80!black}{%
        GeoLayers \api{} $\to$ 7 capas GeoJSON georreferenciadas%
    }}
};

%% ——— Flechas dominio ↔ adaptadores internos ———
\draw[flecha] (api.east)  -- (-4.0,  0.0);
\draw[flecha] ( 4.0, 0.0) -- (goclient.west);
\draw[flecha] ( 4.0, 0.0) -- (geojson.west);
\draw[flecha] ( 4.0, 0.0) -- (cache.west);

%% ——— Flecha visualización (desciende desde el lado inferior del hexágono) ———
\draw[flechaviz] (0, -3.464) -- (viz.north);

%% ——— Flechas externas ———
\draw[flechaext] (cliente.east)  -- (api.west);
\draw[flechaext] (goclient.east) -- (apimetro.west);

%% ——— Borde Docker (solo las tres capas internas) ———
\begin{scope}[on background layer]
    \node[draw=vftVerde!40, dashed, rounded corners=14pt, fill=vftVerde!2,
          fit={(api)(hexbox)(goclient)(geojson)(cache)},
          inner sep=22pt] (docker) {};
\end{scope}
\node[font=\scriptsize\bfseries, color=vftVerde!65, anchor=north]
    at (docker.north) {Docker};

\end{tikzpicture}
}%
    \caption[Arquitectura hexagonal de VFTModel]{Arquitectura hexagonal de
    VFTModel: el hexágono central es el \textbf{Dominio} (\texttt{src/core/}),
    núcleo analítico aislado de cualquier tecnología externa. El
    \textbf{adaptador primario} (izquierda) expone los algoritmos vía \api{}
    \textsc{rest}; los \textbf{adaptadores secundarios} (derecha) conectan el
    núcleo con Apimetro, los archivos locales de respaldo y la capa de caché. La
    zona inferior —fuera del contenedor Docker— representa los
    \textbf{Conectores de Visualización}: \termino{Transport-GIS-MJG}, Tableau y
    \textsc{qgis} consumen los resultados topológicos a través de los siete
    \termino{geo-layers} que expone la \api{}. El nombre \termino{hexagonal}
    proviene de esta disposición en puertos alrededor del núcleo, siguiendo el
    patrón \termino{Ports and Adapters} propuesto por Alistair Cockburn.}
    \label{fig:vft-arquitectura}
    \fuentefigura{Elaboración propia.}
\end{figure}

El beneficio práctico de esta separación es la reproducibilidad: si el formato
de datos de Apimetro cambia, o si los resultados se requieren en un visor
cartográfico distinto, solo se modifica la capa correspondiente sin alterar el
núcleo matemático.

\subsubsection{Flujo de Procesamiento}
\label{subsec:vft-flujo}

El procesamiento de una solicitud de análisis sigue una cadena de seis
transformaciones sucesivas, desde la ingesta de datos geográficos crudos hasta
la entrega de métricas finales, como se esquematiza en la Figura~
\ref{fig:vft-flujo}.

\begin{figure}[H]
    \centering
    \resizebox{\linewidth}{!}{%
    % Diagrama: Flujo de procesamiento de VFTModel (6 etapas encadenadas)
% Usado en: 4-Capas-Codigo/4-2-VFTModel.tex  (fig:vft-flujo)
% Nota: el \resizebox{\linewidth}{!}{...} permanece en el capítulo, envolviendo este \input
\begin{tikzpicture}[
    node distance = 0.5cm and 0.35cm,
    caja/.style = {
        rectangle, draw=unamAzul, fill=unamAzul!7,
        text width=1.95cm, align=center,
        minimum height=1.3cm, rounded corners=4pt,
        font=\small
    },
    flecha/.style = {-Stealth, thick, color=unamAzul!70},
    etiq/.style = {
        font=\scriptsize\itshape, color=gray,
        align=center, text width=1.95cm
    }
]
    \node[caja] (apy) {\textbf{Apimetro}\\{\tiny Go Backend}};
    \node[caja, right=of apy]  (val) {\textbf{Validación}\\{\tiny Pydantic}};
    \node[caja, right=of val]  (grf) {\textbf{Constructor}\\{\tiny del Grafo}};
    \node[caja, right=of grf]  (imp) {\textbf{Impedancia}\\{\tiny Motor VFT}};
    \node[caja, right=of imp]  (ana) {\textbf{Analizadores}\\{\tiny Topológicos}};
    \node[caja, right=of ana]  (res) {\textbf{Resultados}\\{\tiny JSON / CSV}};

    \node[etiq, below=0.18cm of apy]  {Estaciones\\y Líneas\\\geojson{}};
    \node[etiq, below=0.18cm of val]  {Contratos\\de datos};
    \node[etiq, below=0.18cm of grf]  {Grafo\\NetworkX};
    \node[etiq, below=0.18cm of imp]  {Pesos en\\minutos};
    \node[etiq, below=0.18cm of ana]  {8 Indicadores};
    \node[etiq, below=0.18cm of res]  {{\upshape\api{}} REST};

    \draw[flecha] (apy) -- (val);
    \draw[flecha] (val) -- (grf);
    \draw[flecha] (grf) -- (imp);
    \draw[flecha] (imp) -- (ana);
    \draw[flecha] (ana) -- (res);
\end{tikzpicture}
}%
    \caption[Flujo de procesamiento de VFTModel]{Cadena de transformación de
    VFTModel: desde los datos geográficos crudos de Apimetro hasta los
    indicadores topológicos publicados vía \api{} \textsc{rest}.}
    \label{fig:vft-flujo}
    \fuentefigura{Elaboración propia.}
\end{figure}
